\documentclass[UTF8,a4paper,10pt]{ctexart}
\usepackage{../../../../../../common/ipara}
\usepackage{amsmath,amssymb,booktabs,tabularx}
\begin{document}
\section*{H-B03\quad 微分与累积：基本定理及正则性边界}
\textbf{边界：}本桥连接局部变化率与累积量，导数和积分的独立机制分别回到 Topic03、Topic05。
\begin{tcolorbox}[title=三类对象必须分开]
\begin{tabularx}{\linewidth}{@{}lX@{}}
原函数 $F$ & $F'=f$，要求的是逐点导数关系；\\
定积分 $\int_a^b f$ & 一个数，要求区间上的可积性；\\
变上限函数 $G(x)=\int_a^x f(t)\,dt$ & 先由累积定义函数，再讨论其连续性或可导性。
\end{tabularx}
\end{tcolorbox}
\subsection*{1.\ 两个方向}
若 $f$ 在 $[a,b]$ 连续，定义 $G(x)=\int_a^x f(t)\,dt$，则
\[
G'(x)=f(x),\qquad \int_a^b f(x)\,dx=F(b)-F(a)\quad(F'=f).
\]
第一式是变上限积分的局部结论，第二式是 Newton--Leibniz 的累计结论；不能把一个对象的条件自动套给另一个对象。
\subsection*{2.\ 正则性梯}
\[
f\text{ 连续}\Rightarrow G\text{ 可导},\qquad
f\text{ Riemann 可积}\Rightarrow G\text{ 连续},
\]
但“$f$ 可积”不等于“$G'=f$ 处处成立”。例如阶跃函数在跳跃点不可连续，变上限积分在该点通常只能讨论单侧导数或不可导。
\subsection*{3.\ 反向使用}
遇到 $\int f$：先问是否能找到原函数；遇到 $G'(x)$：先识别 $G$ 是否为变上限积分；遇到积分恒等式：再检查端点、参数与可积区间。
\subsection*{4.\ 最小例子与调用}
令 $f(t)=0$（$t<0$），$f(t)=1$（$t\ge0$），并定义 $G(x)=\int_{-1}^{x}f(t)\,dt$。则
\[
G(x)=\begin{cases}0,&x\le0,\\x,&x>0,\end{cases}
\]
$G$ 连续但在 $0$ 不可导，说明“可积 $\Rightarrow G$ 连续”不能升级为“$G'=f$ 处处”。题面混用原函数、定积分和变限函数时调用。
\begin{tcolorbox}[title=检查句]
我当前操作的对象是“一个数、一个函数，还是一个原函数”？所需条件是连续、可积，还是存在逐点导数？
\end{tcolorbox}
\end{document}
